\subsection{Examples of flat modules}

\begin{enumerate}
\def\labelenumi{(\arabic{enumi})}
\item
  For any ring A, \(A_{d}\) is clearly a flat A-module (Algebra, Chapter II, § 3, no. 4, Proposition 4). Then it follows from Proposition 2, (i) of no. 3 that every free right A-module, and more generally every projective right A-module (Algebra, Chapter II, § 2, no. 2), is a flat A-module.
\item
  If A is a semisimple ring (Algebra, Chapter VIII, § 5, no. 1, Definition 1) every right A-module E is semisimple and hence a direct sum of simple modules; as each of these latter is isomorphic to a direct factor of \(\mathbf{A}_d\) (ibid., § 5, no. 1, Proposition 6), E is projective and therefore flat by (1) (cf.~Exercise 16).
\end{enumerate}

*(3) In Chapters II and III we shall study in detail two important examples of flat A-modules: rings of fractions \(S^{-1}A\) and Hausdorff completions \(\hat{A}\) of \(A\) for \(\mathfrak{I}\) -adic topologies.*

PROPOSITION 3. Let A be a ring and E a right A-module.

\begin{enumerate}
\def\labelenumi{(\roman{enumi})}
\item
  Suppose that E is flat. For every element a of A which is not a right divisor of 0(*), the relations \(x \in \mathbf{E}\) , \(xa = 0\) imply \(x = 0\) .
\item
  Suppose that A is an integral domain in which every finitely generated ideal is principal (for example a principal ideal domain (Algebra, Chapter VII, § 1, no. 1)). Then for E to be flat it is necessary and sufficient that E be torsion-free.
\end{enumerate}

We prove (i). Let \(v: \mathbf{A}_s \to \mathbf{A}_s\) be the left A-module homomorphism \(t \mapsto ta\) ; the hypothesis implies that \(v\) is injective. As E is flat, the homomorphism \(1_{\mathbf{E}} \otimes v: \mathbf{E} \otimes_{\mathbf{A}} \mathbf{A}_s \to \mathbf{E} \otimes_{\mathbf{A}} \mathbf{A}_s\) is also injective. When \(\mathbf{E} \otimes_{\mathbf{A}} \mathbf{A}_s\) is canonically identified with E, \(1_{\mathbf{E}} \otimes v\) becomes the endomorphism \(x \mapsto xa\) of E. Thus the relation \(xa = 0\) implies \(x = 0\) .

We prove (ii). By (i), if E is flat, E is torsion-free. Conversely, let E be a torsion-free A-module; we verify that, for every finitely generated ideal a of A, the canonical homomorphism \(E \otimes_{A} a \to E\) is injective (no. 3, Remark 1). This assertion is obvious if \(a = (0)\) ; otherwise, by hypothesis a = Aa for some \(a \in A\) , \(a \neq 0\) , and \(t \mapsto ta\) is then an isomorphism v of A onto a; using i to denote the canonical injection \(a \to A\) , \(i \circ v\) is the homothety with ratio a on E and is injective since E is assumed to be torsion-free. Then \(1_{E} \otimes (i \circ v) = (1_{E} \otimes i) \circ (1_{E} \otimes v)\) ; as \(1_{E} \otimes v\) is an isomorphism, \(1_{E} \otimes i\) is injective, which completes the proof.

Example. Applying Proposition 3 to the ring \(\mathbf{Z}\) , it is seen that \(\mathbf{Q}\) is a flat \(\mathbf{Z}\) -module, but that \(\mathbf{Z}/n\mathbf{Z}\) (for \(n \geqslant 2\) ) is not a flat \(\mathbf{Z}\) -module.

